% Research notes: the serving protocol in full (moved unchanged from the paper's Appendix E, except
% cross-references; owner: appendix editor).
\section{The serving protocol}\label{notes:serving}
Appendix~\paperref{app:serving} of the paper gives the stock serving frontier and what it measures. This section holds the protocol behind it in full: the workload and its splits, the metrics, the launch checks and capacity, the configuration search, the output comparisons against stock references, the GSM8K quality check, the frontiers and the measurement hazards. Every arm in it is stock or exact up to rounding; none uses the certified head.

Every served speedup is to be measured against a server, not a kernel; the only served runs so far are of stock and exact-up-to-rounding arms. This section defines the measurement: what runs, on which prompts, how the latency--throughput frontier is computed, how the baselines are made strong before anything is compared with them, and which quality check a change must pass.

\paragraph{Setup and workload}\label{sec:serving-setup} All runs use the machine and engine of Section~\paperref{sec:setup} of the paper, serving the BF16 target at its pinned revision with FlashInfer attention~\cite{flashinfer}; the vision encoder, which text requests never reach, uses Triton attention. The client is AIPerf~\cite{aiperf}. Speculative acceptance depends on what the model is writing, so random-token prompts would measure the wrong thing. The benchmark uses a frozen set of 1,857 public prompts in three domains: chat (the first turns of the 80 MT-Bench questions~\cite{mtbench} and English root prompts from OASST1~\cite{oasst1}), code (the 164 HumanEval problems~\cite{humaneval} and MBPP test tasks with their assertions~\cite{mbpp}) and mathematics (GSM8K training questions with a request to box the final answer~\cite{gsm8k}; the GSM8K test split is reserved for the quality check, so no workload prompt is a quality problem). Each source is pinned to a Hugging Face revision and the builder reproduces the files byte for byte. Each domain is cut into disjoint warm-up (43 prompts), tuning (192) and confirmation (384) sets, and every split interleaves the domains so that any prefix of it is balanced, which matters because low-concurrency points use only a prefix. Configurations are chosen on the tuning split and reported on the confirmation split. Templated prompts are short: the confirmation split has a mean of 87.9 prompt tokens (median 78, 99th percentile 283), with chat shorter (mean 47) than code (130) and mathematics (86)~\ev{workload}.

Requests use the model's chat template with thinking enabled (the template's default) and no reasoning parser, so the thinking text is streamed and counted as output. Decoding is greedy. The throughput panel fixes the output at 512 tokens with \code{ignore\_eos}, and every point checks that each response has exactly that length. A fixed length is a faithful sample of the model's behaviour only while the model has not chosen to stop, so a natural-stopping run on 576 tuning prompts measured where greedy outputs end. Only 3.8\% stopped before 512 tokens and 0.7\% had entered a detected repetition loop by then (a 256-token window with fewer than 30\% distinct 8-grams), so almost all of the fixed-length panel is text the model would have produced anyway. The same run shows why greedy decoding is not used for the quality check: 142 of the 576 outputs (25\%; 36\% in mathematics, 32\% in chat, 6\% in code) ran to the 16,384-token limit, a degeneration Qwen's model card warns about for greedy thinking~\cite{qwencard}. The run used an earlier version of the workload, whose mathematics prompts came from the GSM8K test split, and other processes used up to 24 CPU cores during it, so it is a length measurement, not a timing~\ev{natural-lengths}.

\paragraph{Metrics}\label{sec:serving-metrics} For request $i$ at a sweep point let $n_i$ be the number of generated tokens the server reports in the \code{completion\_tokens} field of its usage record, $s_i$ the time the client sends it, $f_i$ the arrival of its first non-empty streamed chunk and $e_i$ the interval from $s_i$ to its last non-empty chunk. Over the $N$ measured requests the frontier plots
\begin{equation}
x=\frac1N\sum_{i=1}^N\frac{n_i}{e_i},\qquad y=\frac{\sum_in_i}{\max_i(s_i+e_i)-\min_is_i},\label{eq:frontier}
\end{equation}
so $x$ is the mean per-request output rate \emph{including} the time to first token, and $y$ is output tokens per second on the one GPU. Two companions are reported and never substituted for them. The decode rate $x_{\rm dec}=\frac1N\sum_i(n_i-c_i)/(e_i-(f_i-s_i))$, where $c_i$ is the exact token count of the first chunk from per-chunk usage (a speculative step streams several tokens per chunk), excludes queueing and prefill; one token per chunk is never assumed. The steady-state throughput counts only the tokens streamed between the first and the last send, when the client holds its full concurrency, and so excludes the drain at the end of a point. Each point also records time-to-first-token, inter-token and end-to-end latency percentiles, failures, the per-request speculative statistics the server attaches to each stream (verify steps, accepted drafts, and accepted length as output tokens per verify step with the bonus token included), and deltas of the server's counters, among them the fraction of decode passes that replayed a CUDA graph. The scheduler's logged generation rate while the batch is nearly full is recorded as a diagnostic of what the GPU sustains; clients do not see it, and it is never a frontier metric.

\paragraph{Protocol and baselines}\label{app:serving-protocol}\label{sec:serving-protocol} A server arm is a model revision plus a declarative set of launch flags. The harness starts it once, verifies it, and sweeps client concurrency against it in closed loop; the confirmation sweep covers client concurrencies 1, 2, 4, 8, 16, 32, 48, 64, 96 and 128, each arm over the concurrencies where it competes (\path{bench/campaigns/confirm.sh}). Before trusting a launch it checks that decode CUDA graphs were captured for every batch size up to the server's capacity (the target decode graph, or the target-verify, draft-decode and draft-extend graphs under speculation), that prefill graphs were captured, that the overlap scheduler is running (the scheduler reports it, no fallback is logged, and a stack sample of the scheduler process shows its overlap loop), and that the recurrent-state cache did not reduce the capacity.

Capacity needs care on this model. The attention cache is not the constraint: 32~KiB per token, and the default memory split leaves room for about a million tokens. The recurrent state is. With the prefix cache and the overlap scheduler on, SGLang reserves five FP32 state slots per running request (three for prefix-cache retention and two ping-pong buffers), while decode touches one; the 667 slots the default split gives the state pool therefore cap the server at 133 running requests (code reading, confirmed against the startup log)~\ev{profile-attr}. Speculative verification adds one intermediate state per draft token, and SGLang resets speculative servers to 48 running requests, and logs it, unless the capacity is given; the capacity each arm would reach at default settings has not been measured. Feasibility probes show what launches: plain decoding serves 256 concurrent requests with the prefix cache and 1,024 without it, MTP speculation serves 256 without the prefix cache, and an MTP server with BF16 state at 512 fails at verify-graph capture on a FlashInfer workspace overflow. These probes establish what runs, not performance~\ev{probes}. Every arm is therefore launched at a capacity of 128 running requests, the top of the client sweep (64 for one DFlash arm, below), with the state cache sized for it and the remaining static memory given to the attention cache and the verify states.

Every arm streams every fourth token (\code{stream-interval = 4} in \code{bench/arms.toml}). The server still sends the first token at once and the last at finish, so the times that Equation~\eqref{eq:frontier} uses do not depend on the interval, and the choice takes load off the single-threaded streaming path (below). Each point sends $\max(64,8c)$ measured requests, the first prompts of the confirmation split in order, after one wave of warm-up requests at the same concurrency, and the prefix cache is flushed before every point so that no measured prompt is served from cache. The confirmation sweep runs in three sessions. Each session launches every tuned arm afresh in the same exclusive hold as its matched plain baseline (the FlashInfer plain arm for FlashInfer arms, the Triton one for Triton arms), in reversed order in alternate sessions, so that drift does not align with an arm. Within a session each arm runs once, over its concurrencies in ascending order; the harness alternates the order of concurrencies only between repeats within one launch, which the confirmation sweep does not use. Speedups are ratios within a session; a pair with an invalid point on either side is dropped and counted, never replaced by another run, and points at which the attention cache limited the batch are flagged from the scheduler log. Timed runs hold the GPU exclusively, and because the scheduler, tokenizer manager and client are single-threaded CPU loops, each point also waits until other processes use fewer than two cores and records the foreign load it saw.

The speedup denominator is the strongest existing configuration, not a default. Arm~A is ordinary autoregressive decoding. Arm~B is native MTP speculation, whose draft steps reuse the target's tied head. Arm~C is the public DFlash drafter (Section~\paperref{sec:setup} of the paper). The search ran on the tuning split~\ev{tuning}: 23 configurations, each launched once and measured at client concurrency 1, 8, 32 and 128 (64 instead of 128 for the two block-16 configurations) with $\max(32,4c)$ requests per point and 512 output tokens. It covered MTP chain depth (two to five steps with the prefix cache, three, five and seven without it), the prefix cache, SGLang's replayed GDN verification and replayed decoding, its batch-adaptive depth, Triton against FlashInfer attention, DFlash blocks of 4, 8 and 16 tokens and FA4 attention for the DFlash drafter. A one-step chain is not in the search, because it failed a launch check that required a draft-decode graph such a chain does not capture; replayed verification was refused at launch for DFlash and with FlashInfer's GDN decode kernel. Draft trees (top-$k$ 2 and 4, which run on this hybrid model, with accepted lengths of 3.45 and 3.56 at concurrency 1 in the feasibility probes~\ev{probes}), the verify attention mode, the prefill chunk size and graph backend were not searched. With one launch per configuration, Table~\ref{tab:tuning} is selection evidence, not a result.

\begin{table}[!t]
\small\centering
\caption{Configuration search on the tuning split: output throughput $y$ (Equation~\eqref{eq:frontier}) from one launch per configuration with $\max(32,4c)$ requests per point, and the accepted length at concurrency 32. FlashInfer attention unless Triton is named; MTP chains draft one token per step; RV: replayed GDN verification. ``excl.'': excluded because other processes used more than two CPU cores on average, by a sampler that then overcounted, so these may be false positives. $^\dagger$Limited by the attention cache, which without a cap had shrunk to 54K tokens as the state buffers grew: at most 123 requests ran, with retractions. $^\ddagger$Capacity 64; the last throughput column is concurrency 64. Selection evidence, not results~\ev{tuning}.}
\label{tab:tuning}
\vskip 0.05in
\begin{tabular}{@{}llcrrrrc@{}}\toprule
& & & \multicolumn{4}{c}{$y$ (tokens/s) at concurrency $c$} & Accepted\\\cmidrule(lr){4-7}
Arm & Configuration & Prefix cache & 1 & 8 & 32 & 128 & length\\\midrule
Plain & FlashInfer & on & 282 & 1{,}982 & 6{,}087 & 13{,}421 & --\\
& FlashInfer & off & 282 & 2{,}004 & excl. & 13{,}844 & --\\
& Triton & off & 286 & 2{,}018 & 6{,}241 & 13{,}846 & --\\
& FlashInfer, replayed decoding & off & 244 & 1{,}807 & 5{,}909 & 14{,}981 & --\\\addlinespace
MTP & 2 steps & on & 385 & 2{,}288 & 5{,}455 & 9{,}153 & 2.67\\
& 3 steps & on & 461 & 2{,}482 & 5{,}839 & 9{,}593 & 3.28\\
& 4 steps & on & 465 & 2{,}608 & 5{,}761 & 9{,}276 & 3.72\\
& 5 steps & on & 468 & 2{,}575 & 5{,}684 & 8{,}719$^\dagger$ & 4.11\\
& 3 steps & off & 466 & 2{,}534 & 6{,}127 & 10{,}106 & 3.27\\
& 5 steps & off & 481 & 2{,}699 & 6{,}087 & 9{,}456 & 4.14\\
& 7 steps & off & 465 & 2{,}514 & 5{,}690 & 8{,}493 & 4.62\\
& 3 steps, RV & on & 446 & 2{,}646 & 6{,}469 & 11{,}941 & 3.28\\
& 5 steps, RV & on & 462 & 2{,}666 & 6{,}502 & 11{,}186 & 4.11\\
& 3 steps, RV & off & 456 & 2{,}667 & 6{,}925 & 13{,}011 & 3.29\\
& 3 steps, RV, Triton & off & 535 & 3{,}093 & 7{,}103 & 11{,}353 & 3.27\\
& 4 steps, RV, Triton & off & 552 & 3{,}104 & 6{,}958 & 10{,}804 & 3.73\\
& Adaptive depth, RV & off & 474 & 2{,}603 & 5{,}968 & 10{,}010 & 1.97\\\addlinespace
DFlash & Block 4 & off & 498 & 2{,}682 & 6{,}380 & 11{,}477 & 3.28\\
& Block 8 & off & 686 & 3{,}443 & 6{,}727 & 10{,}226 & 4.75\\
& Block 8, FA4 draft attention & off & 769 & excl. & 6{,}983 & 10{,}594 & 4.75\\
& Block 8, Triton & off & 807 & 3{,}613 & 6{,}307 & 8{,}957 & 4.75\\
& Block 16$^\ddagger$ & off & excl. & 3{,}398 & 5{,}386 & 6{,}737 & 5.71\\
& Block 16, Triton$^\ddagger$ & off & 851 & 3{,}639 & 5{,}199 & 6{,}397 & 5.72\\\bottomrule
\end{tabular}
\end{table}

Three regularities decide the arms. First, speculation pays at low concurrency and loses at high. MTP with three steps delivers 1.64 times plain decoding's $y$ at concurrency 1 (461 against 282 tokens/s); with snapshot verification it falls behind from concurrency 32 (5,839 against 6,087), with replayed verification and no prefix cache it still leads there (6,925) and trails at 128 (13,011 against 13,844), and at concurrency 128 every speculative configuration trails plain decoding. Second, replayed verification, which removes the per-draft-token state snapshots, recovers 24\% at concurrency 128 for three steps (11,941 against 9,593 tokens/s) at a 3\% cost at concurrency 1, and depth 3 beat depths 4 and 5 at concurrency 128 in every pair the attention cache did not limit, while trailing them by at most 4\% at concurrency 1 and 7\% at 8; in the low-concurrency arm's own configuration (replayed verification, Triton), depth 4 led depth 3 by 3.2\% at concurrency 1 and 0.4\% at 8 and trailed it by 2\% at 32, so one depth serves both MTP arms. Third, Triton attention helps every speculative family at low concurrency and hurts at high concurrency, and plain decoding is indifferent to it. DFlash at block 16 under Triton has the highest throughput of any configuration at concurrency 1 and 8, and at concurrency 1 a per-request rate 3.35 times that of plain decoding under Triton (957 against 286 tokens/s in $x$); among DFlash configurations block 4 leads at concurrency 128 (11,477 tokens/s, 8\% above block 8 with FA4 draft attention).

The chosen arms (\path{bench/arms.toml}) all run without the prefix cache, which this workload does not reuse and whose state slots deeper drafts and replayed verification need. Arm~A is plain decoding under FlashInfer, with a Triton twin as the matched baseline for Triton arms. Arm~B is MTP with three steps and replayed verification, under FlashInfer for high concurrency and under Triton for low. Arm~C is DFlash at block 8 with FA4 draft attention for every concurrency, at block 16 under Triton for low concurrency, capped at 64 running requests because its 16 verify states per request take about 0.8~GB, and at block 4 for high concurrency. Two of these choices change the target's arithmetic: replayed verification computes its verify outputs with a chunked transform on TF32 tensor cores (code reading at the pinned commit), and Triton attention for the target replaces the target's attention kernel. Both are exact up to rounding by the comparison below, as is replayed decoding, so every confirmed arm is stock or exact up to rounding. MTP with SGLang's snapshot verification under FlashInfer is kept as a stock reference, and plain decoding with replayed decoding, the fastest configuration at concurrency 128 on the tuning split and the slowest at 1, as a variant.

\paragraph{Output equality}\label{sec:serving-equality} Each arm that changes the target's arithmetic was compared with a matched stock reference on the 320 divergence prompts of Appendix~\paperref{app:divergence} of the paper, decoded greedily for 256 tokens with the top five log-probabilities at concurrency 1 and without the prefix cache, like every arm: plain levers against stock plain decoding, buffered MTP against stock MTP with three steps, and Triton DFlash at block 16 against stock DFlash at block 16, both DFlash runs at a capacity of 4~\ev{bench-equality}. The class rule was fixed after the first comparisons had been seen. It rests on the divergence classes of Appendix~\paperref{app:divergence} of the paper, which predate them: an arm is exact up to rounding if every prompt's first divergence is at an exact tie, within one BF16 step, or a near tie with both margins within two steps, and no output differs in length; it is lossy otherwise. Every arm passes, with no event above a near tie (Table~\ref{tab:equality}). The table also gives each rate's ratio to the batch-shape floor of 3.42 per 1,000 (plain decoding at concurrency 1 against 32 on one server with the prefix cache on and at most 16 running requests). Every ratio's 95\% interval includes 1, and the ratios inform the class rather than decide it. The ratios depend on that floor's cap: with at most 8 running requests and fixed pools the same pair diverges at 0.63 per 1,000 (Appendix~\paperref{app:divergence} of the paper), while the class rule does not use the floor's rate. The classes above were measured with servers capped at 16 running requests (4 for DFlash), and none of these arms has been compared at a cap of 8; only stock plain decoding and stock MTP were rerun there, and they too diverge only at rounding-level events, an exact tie, one BF16 step or two steps at most~\ev{state-pinned}. Replayed plain decoding was first classed lossy against the plain run of Appendix~\paperref{app:divergence} of the paper, which has the prefix cache on: one prompt's first divergence was large. At that position the prefix-cache-on run's own token rests on a logit gap of 0.0625, which stock plain decoding without the prefix cache resolves the other way, and with the prefix cache on a request's log-probabilities can depend on earlier requests (Appendix~\paperref{app:divergence} of the paper). Against the reference without the prefix cache, replayed plain decoding first differs from it in that prompt at a tie further on, and its rate is the floor's. The two stock plain references, with and without the prefix cache, differ in 96 of 320 prompts at 1.69 per 1,000, all rounding-level.

\begin{table}[htbp]
\small\centering
\caption{Greedy outputs against matched stock references on the 320 divergence prompts (256 tokens, top-five log-probabilities, concurrency 1, no prefix cache; servers capped at 16 running requests, 4 for DFlash): divergences per 1,000 compared tokens and the ratio to the batch-shape floor of 3.42 per 1,000, with its 95\% interval where computed. Every arm that changes the target's arithmetic is exact up to rounding by the class rule. Measured~\ev{bench-equality}.}\label{tab:equality}
\vskip 0.05in
\begin{tabular}{@{}llrr@{}}\toprule
Arm & Reference & Per 1,000 & Ratio to the floor\\\midrule
\multicolumn{4}{@{}l}{\emph{Arms that change the target's arithmetic}}\\
MTP, three steps, buffered verification & stock MTP, three steps & 3.79 & 1.11\\
the same with Triton attention & stock MTP, three steps & 4.05 & 1.19\\
DFlash, block 16, Triton attention & stock DFlash, block 16 & 3.82 & 1.12\\
Plain decoding, Triton attention & stock plain decoding & 3.89 & 1.14\\
Plain decoding, replayed & stock plain decoding & 3.40 & 0.99\\\addlinespace
\multicolumn{4}{@{}l}{\emph{Stock speculation}}\\
MTP, three steps & stock plain decoding & 3.95 & 1.15 (0.93--1.43)\\
DFlash, block 16 & stock plain decoding & 3.92 & 1.15 (0.93--1.42)\\\bottomrule
\end{tabular}
\end{table}

\paragraph{Speculation against plain decoding} Stock speculation against the same plain reference diverges at 1.15 times the floor's rate both for MTP with three steps and for DFlash at block 16 (Table~\ref{tab:equality}). Against the prefix-cache-on run the ratios were 1.29 and 1.28, with intervals above 1, so that apparent excess came at least partly from the reference's prefix-cache setting, not from speculation: at the pinned commit that setting changes the recurrent prefill of prompts longer than one 64-token chunk, from the first output token, in all 6 such prompts of a 40-prompt check of plain decoding at batch 1 and in 192 of the 193 among the 320 for MTP with three steps at batch 1 with fixed pools (Appendix~\paperref{app:divergence} of the paper). Turning the overlap scheduler off, which at the pinned commit also changes the prefix cache's recurrent-state strategy, changes the same first tokens for MTP. For plain decoding on the same 320 prompts either switch changes the first output token in 192 of the 193 prompts longer than 64 tokens and in none of the 127 shorter ones, as for MTP~\ev{state-pinned}. Against this floor, then, an excess of speculation is not detected, though one of up to about 40\% is not excluded. Three limits apply to this check: the floor is a prefix-cache-on pair on one server, and a floor without the prefix cache was not run; the servers' pools were sized from free memory, not pinned, although for MTP with three steps at concurrency 1 without the prefix cache the run of this comparison with pools sized at start-up and one with fixed pools at a cap of 8 agree bitwise on all 320 prompts~\ev{state-pinned}; and the DFlash runs used a capacity of 4 against plain decoding's 16. Whether speculation shows an excess depends on the floor, however. With fixed pools and at most 8 running requests, both runs with the prefix cache on, plain decoding at concurrency 1 against 32 diverges at 0.63 per 1,000 while MTP with three steps against plain decoding at concurrency 1 diverges at 3.68, about six times that floor (derived), and every MTP configuration diverges from plain decoding at 3.5 to 4.0 per 1,000, about as often as it diverges from itself between concurrency 1 and 32 (3.1 to 3.5)~\ev{state-pinned}. The speculative rate is about the same in both regimes, and what moves is plain decoding's own batch-shape divergence.

\paragraph{Quality check}\label{sec:serving-quality} Output length and throughput say nothing about whether a change preserved the model. The quality check is the full GSM8K test split (1,319 problems) scored with the GSM8K grader of the \code{sgl-eval} package, version 0.1.2 (a zero-shot prompt, the last boxed answer, symbolic equality) with thinking on, natural stopping at 16,384 tokens, and Qwen's recommended thinking-mode sampling for precise tasks (temperature 0.6, top-$p$ 0.95, top-$k$ 20) under a fixed request seed. Two runs are compared problem by problem: the accuracy difference, an exact McNemar test on the discordant problems, and the share of identical final answers. Sampled accuracy tests equality in distribution for a lossless change and prices a lossy one; it cannot show that greedy outputs are unchanged, which the output comparisons of Appendix~\paperref{app:divergence} of the paper address. No arm differs detectably from plain decoding~\ev{bench-quality}. Two launches of plain decoding score 89.99\% and 90.30\%, buffered MTP 89.54\%, MTP with snapshot verification 88.93\%, DFlash at block 8 89.69\% and replayed plain decoding 91.05\% (run while it was still classed lossy). Every paired difference against either plain run lies between $-1.36$ and $+1.06$ points, with exact McNemar $p\ge0.18$. A non-significant test excludes nothing by itself. The paired 95\% intervals (Wald intervals on the discordant problems, derived from the committed counts) have lower bounds between $-0.7$ and $-3.3$ points, the lowest for MTP with snapshot verification against the second plain run ($-1.36$ points, interval $-3.27$ to $+0.54$), so the data exclude losses larger than about 3.3 points for every arm and cannot exclude smaller ones. Two sampled runs of one arm agree on the final answer for only 77\% of problems, so the fixed request seed does not make these runs reproducible, and 17 to 19\% of outputs in every arm reach the 16,384-token limit.

\paragraph{Frontiers}\label{sec:serving-frontiers} The baseline frontier is the upper envelope over arms A, B and C, since no arm need be best at every concurrency: speculation saves weight traffic, which matters most at low concurrency, and adds per-request state traffic, which matters most at high concurrency (Table~\paperref{tab:bytes} of the paper). Every arm carries an exactness class in \path{bench/arms.toml}: \emph{stock} when it sets only flags known to leave the target's arithmetic as in plain decoding with FlashInfer target attention and stock state handling, \emph{exact up to rounding} by the comparison above, \emph{pending} until compared, and \emph{lossy} otherwise. All nine confirmed arms are stock or exact up to rounding, so the envelope over exact arms equals the envelope over all arms. The frontier over three confirmation sessions, with paired ratios against the matched plain baselines and the run-to-run variance, is committed~\ev{bench-frontier}; only points with three valid sessions rank, so replayed plain decoding, with one session, does not rank yet; its two further sessions are pending~\pend{bench-replayssm}. No served result for the certified head exists yet: its before-and-after frontiers in MTP drafting, verification and plain decoding, with output-equality checks at the same batch shape, have not been run~\pend{bench-cert}.



\paragraph{Measurement hazards}\label{app:serving-hazards}\label{sec:serving-hazards} Two effects outside the GPU shaped this protocol. At high concurrency plain decoding can be limited by the streaming path rather than the GPU; speculative arms stream several tokens per chunk and are less exposed, so a frontend bottleneck would inflate their measured advantage. A diagnosis with one plain server at client concurrency 256 on a quiet host measured the effect~\ev{frontend}. With one-token chunks the tokenizer-manager process (HTTP and the OpenAI serving layer) ran at a full core and the client received about 85\% of the server's full-batch decode rate; streaming every four tokens brought that process to 0.35 cores and gave the same throughput as a non-streaming client, about 92\% of the full-batch rate, the rest being the prefill of each new wave and the drain, which $y$ includes by definition. More tokenizer and detokenizer workers made it worse, and more client workers changed nothing. Every row of that diagnosis is a single point; concurrency above 256 on a quiet host has not been measured~\pend{bench-frontend}. The same diagnosis attributes the low throughput of the earlier feasibility probes at concurrency 256 to 1,024 to CPU contention on the host, not to a frontend limit. That is the second effect: exclusive use of the GPU does not keep CPU-heavy jobs off the host, which is why every point records foreign CPU load and waits for a quiet machine.
